% Chapter 4 — Experimental Methodology

\chapter{Experimental Methodology}
\label{chap:methodology}

This chapter describes the experimental setup used in
Chapter~\ref{chap:results}. The experiments combine two one-player reward
processes with repeated play in RPS and RPSLS. Full-information and bandit
learners are evaluated with the same expected action-regret metrics introduced
in Chapter~\ref{chap:background}. Repeated-game experiments additionally track
the empirical joint-action distribution and its distance to the CE and CCE
sets.

\section{Experimental Setup}
\label{sec:experimental-design}

The main trajectory experiments use horizon \(T=10^6\), 12 independent
replicates, and base seed 42. All random-number streams are deterministically
derived from this base seed, so rerunning the same experiment configuration
reproduces the same stochastic trajectories. The one-player experiments use
\(K=10\) actions unless the action-space size is varied. The full-information
and bandit algorithm sets are those summarized in
Chapter~\ref{chap:algorithms}.

The additional scaling experiments are:

\begin{itemize}
    \item \textbf{Action-space scaling:} HFA with
    \(K\in\{2,3,5,10,20,30,50\}\) and \(T=10^6\).
    \item \textbf{One-player horizon scaling:} HFA and LRW with \(K=10\).
    \item \textbf{Repeated-game horizon scaling:} RPSLS self-play and selected
    cross-play profiles.
\end{itemize}

The horizon-scaling experiments use
\[
\{1000,2000,3000,5000,10^4,2\!\cdot\!10^4,3\!\cdot\!10^4,
5\!\cdot\!10^4,10^5,2\!\cdot\!10^5,3\!\cdot\!10^5,
5\!\cdot\!10^5,10^6\}.
\]
Each horizon is run as an independent experiment: the learner is reinitialized,
and any fixed-horizon parameter is computed from that run's own \(T\).

\section{One-Player Environments}
\label{sec:experimental-environments}

The two one-player environments are deliberately different in how the reward
sequence is generated.

\subsection{Adaptive Historical-Frequency Environment}

The historical-frequency adversary (HFA) is \emph{adaptive but
non-anticipating}. Before round \(t\), it ranks the actions by how often they
have been selected in the previous rounds. The \(\lceil K/2\rceil\) most
frequently selected actions receive reward \(0\), while the remaining actions
receive reward \(1\). Ties are resolved by a deterministic rotating
order.

Thus \(r_t\) can depend on the learner's past actions, but not on the fresh
random draw \(I_t\) from the current round. For the main \(K=10\) experiments,
five actions are punished each round. HFA is used both for the main one-player
comparison and for the action-space study.

\subsection{Oblivious Lazy Random Walk Environment}

The lazy random walk (LRW) provides an \emph{oblivious} contrast: its reward
sequence is generated independently of the learner's actions. Each action has
its own process on
\[
    \{0,0.1,\ldots,0.9,1\},
\]
starting at \(0.5\). At an interior state, the next reward moves down by
\(0.1\), stays unchanged, or moves up by \(0.1\), each with probability
\(1/3\). At either boundary, it stays or moves one step inward, each with
probability \(1/2\).

For a fixed replicate, the complete reward paths are generated independently
of the learner. Algorithms compared within that replicate therefore face the
same underlying LRW realization. In contrast to HFA, the learner's actions do
not affect future rewards.

\section{Repeated Games}
\label{sec:repeated-game-method}

The repeated-game experiments use Rock--Paper--Scissors (RPS) and
Rock--Paper--Scissors--Lizard--Spock (RPSLS). Raw win/tie/loss payoffs
\(1,0,-1\) are normalized to \(1,0.5,0\), so both games are symmetric
constant-sum games with total payoff \(1\) at every action profile.

For RPS, with actions ordered Rock, Paper, Scissors, player~0's payoff matrix
is
\[
U_{\mathrm{RPS}}
=
\begin{pmatrix}
0.5 & 0   & 1 \\
1   & 0.5 & 0 \\
0   & 1   & 0.5
\end{pmatrix}.
\]
Player~1 receives \(1-U_{\mathrm{RPS}}\) entrywise.

For RPSLS, with actions ordered Rock, Paper, Scissors, Lizard, Spock,
player~0's payoff matrix is
\[
U_{\mathrm{RPSLS}}
=
\begin{pmatrix}
0.5 & 0   & 1   & 1   & 0 \\
1   & 0.5 & 0   & 0   & 1 \\
0   & 1   & 0.5 & 1   & 0 \\
0   & 1   & 0   & 0.5 & 1 \\
1   & 0   & 1   & 0   & 0.5
\end{pmatrix},
\]
and again player~1 receives the complementary payoff \(1-U_{\mathrm{RPSLS}}\).

At each round, both players sample their actions before feedback is produced.
Under full information, a player observes the payoff of every one of its
actions against the opponent's realized action. Under bandit feedback, it
observes only the payoff of its played action. The counterfactual payoff vector
is still retained by the simulator for offline regret evaluation and is never
shown to a bandit learner.

RPS is mainly used to inspect the structure of the empirical joint-action
distribution. RPSLS is the main game for self-play, horizon-scaling, and
cross-play comparisons.

\section{Evaluation Metrics}
\label{sec:evaluation-metrics}

\subsection{Expected Action Regret}

Following the terminology of \citet{greenwald_bounds_2006}, the evaluator
computes \emph{action regret} from the actions actually sampled in each
replicate. It maintains the replacement gains
\[
    G_T(i,j)
    =
    \sum_{t=1}^{T}
    \mathbf 1\{I_t=i\}
    \bigl(r_t(j)-r_t(i)\bigr).
\]
External, internal, and swap action regret are then
\[
    R_T^{\mathrm{ext}}
    =\max_j\sum_i G_T(i,j),
    \qquad
    R_T^{\mathrm{int}}
    =\max_{i,j}G_T(i,j),
\]
and
\[
    R_T^{\mathrm{swap}}
    =\sum_i\max_jG_T(i,j).
\]

The reported trajectories are arithmetic means over the 12 replicates and are
therefore Monte Carlo estimates of the corresponding expected action regret,
\[
    \mathcal R_T^q=\mathbb E[R_T^q],
    \qquad
    q\in\{\mathrm{ext},\mathrm{int},\mathrm{swap}\}.
\]
For swap regret, this placement of the expectation follows the expected swap
regret notion distinguished from pseudo-regret by
\citet{huang_near-optimal_2023}. The main plots use \(R_t/t\); selected
supplementary plots use
\(R_t/\sqrt t\). The same evaluator is used for all algorithms under both
feedback models.

\subsection{Equilibrium Distance}

For repeated games, the empirical joint distribution after \(t\) rounds is
\[
    \widehat{\mu}_t(a)
    =
    \frac{1}{t}\sum_{s=1}^{t}\mathbf 1\{a_s=a\}.
\]
The experiments compute its full-dimensional \(L_1\) distance to the complete
CE and CCE sets:
\[
    d_E(\widehat{\mu}_t)
    =
    \min_{\mu\in E}
    \|\widehat{\mu}_t-\mu\|_1,
    \qquad
    E\in\{\mathrm{CE},\mathrm{CCE}\}.
\]
The corresponding linear programs are solved with HiGHS. Distances are
computed separately for each replicate and then averaged. Joint-action
heatmaps instead average the final empirical distributions across replicates.

\section{Finite-Horizon Scaling}
\label{sec:scaling-and-reporting}

For each independently configured horizon, the final action regret is first
averaged across replicates. Positive replicate-mean endpoints are fitted by
ordinary least squares in log--log coordinates,
\[
    \log R_T=\alpha\log T+\log c,
    \qquad\text{equivalently}\qquad
    R_T\approx cT^\alpha.
\]
The fitted \(\alpha\) and \(c\) are used only as finite-horizon empirical
descriptors, not as estimates of asymptotic regret rates.

Trajectory figures show replicate means without confidence bands. Endpoint
information pages additionally report the sample standard deviation and
standard error. These statistics are descriptive and are not used for formal
hypothesis testing.
